\documentclass[11pt]{article}
\usepackage[margin=1in]{geometry}
\usepackage{amsmath,amssymb,amsthm,hyperref}
\newtheorem{theorem}{Theorem}
\newcommand{\Ph}{\Phi_\ell}
\newcommand{\supp}{\operatorname{supp}}
\title{Computing the second digit from endpoints and testing the cost on the actual hashes}
\author{Project research note}
\date{17 November 2026}
\begin{document}
\maketitle

Let $G$ be a finite abelian group. In its integral group algebra let
$[P][Q]=[P+Q]$, and write $h(T)$ for a coefficient. The operation of
interest at a prime $\ell$ is
\[
   D_\ell(h)=\frac{h^\ell-\Ph(h)}{\ell}\pmod\ell,\qquad
   \Ph\Big(\sum_P h(P)[P]\Big)=\sum_P h(P)[\ell P].
\]
This digit can be computed from endpoints if $h$ is supplied as a
sum of runs along one group step. We give a direct algorithm and a
probability bound showing why applying it to the project's independent
abscissa hashes fails to improve the exponent. We also simplify the
previous prime-square reconstruction before any expensive digit is
computed. Its padding differences eliminate $D_\ell(h)$ itself.

\section{An endpoint formula in characteristic \(\ell\)}
Suppose
\[
 h=\sum_{i=1}^k w_i\sum_{j=0}^{L_i-1}[A_i+jQ],\qquad
 L=\max_i L_i .
\]
Weights are integers, overlaps count with multiplicity, and a length-one
run includes scalar padding at the identity. Empty runs can be dropped.
This representation makes sense for any $Q$, including points of small
order. Use a formal variable $z$ over the group ring and form
\[
 H(z)=\sum_i w_i[A_i]\sum_{j=0}^{L_i-1}z^j,\qquad
 N(z)=(1-z)H(z)=\sum_i w_i[A_i](1-z^{L_i}).
\]
Substituting $z=[Q]$ recovers $h$. Extend $\Ph$ by $\Ph(z)=z^\ell$.
The digit commutes with substitution, by taking the quotient in the
torsion-free integral group algebra before reduction.

The product rule follows by writing $F^\ell=\Ph(F)+\ell D_\ell(F)$
with integer lifts and multiplying two such identities. It gives
\[
 D_\ell(H)=
 \frac{D_\ell(N)-\Ph(H)D_\ell(1-z)}{1-z^\ell}
                  \quad\text{over }\mathbb F_\ell[G][[z]] .
\]
This is the Frobenius-lift product rule with the sign convention in
the definition above. See
\href{https://maths-people.anu.edu.au/~borger/papers/07/forms.pdf}
{Borger and Buium, Differential forms on arithmetic jet spaces}.
The denominator has constant term one in the formal variable. We
perform the division before the substitution which could make the
denominator a zero divisor. The answer is a polynomial of degree at
most $M=\ell(L-1)$.

Combine coinciding terms of $N$ as $c_u[B_u]z^{e_u}$.
For a population $(n_u)$ summing to $\ell$, a pure population (one
$n_u=\ell$) contributes $(c_u^\ell-c_u)/\ell$ times
$[\ell B_u]z^{\ell e_u}$. Every other population contributes
\[
    (\ell-1)!\prod_u\frac{c_u^{n_u}}{n_u!}\,
    \Big[\sum_u n_uB_u\Big]z^{\sum_u n_u e_u} \pmod\ell .
\]
All these factorial denominators are units. There are at most $2k$
endpoint types. The population sum involves endpoints of the runs,
independently of their lengths. For a term $c[B]z^e$, dividing by
$1-z^\ell$ contributes $c$ at $T$ once for each integer $a\ge0$
such that $e+\ell a\le M$ and $B+(e+\ell a)Q=T$. Scanning this
single progression costs $O(L)$ group operations and constant many
working points. It correctly includes every wrap around a short group
cycle. No group index for $T$ is required.

For the correction, the coefficient of $z^j$ in $D_\ell(1-z)$ is
$(-1)^j\binom{\ell}{j}/\ell$ for $1\le j<\ell$, and the coefficient
at $j=\ell$ is $(({-1})^\ell+1)/\ell$. The latter matters at $\ell=2$.
The $i$th run of $\Ph(H)$ then contributes with weight
$w_i\min(L_i,a+1)$ at degree $j+\ell a$, based at $\ell A_i$,
before applying the overall minus sign. This follows by counting the
pairs $v+w=a$ with $0\le v<L_i$, $w\ge0$. Thus this correction uses
only single progressions as well.

\begin{theorem}
Given the run description, the digit at one target costs at most
\[
    (1+L)\left\{\binom{2k+\ell-1}{\ell}+k\ell\right\}
             \operatorname{poly}(\ell,k,\log(2+L),b,\log(2+\max_i|w_i|))
\]
bit operations when group operations on $b$-bit points take polynomial
time. The evaluation can stream the populations and needs no table of
partial group sums. For singleton runs it applies to arbitrary explicit
masks. For $k=O(m)$, $\ell=O(m)$ and $L=O(\tau)$ the logarithm of
the cost is $b/m+O(m+\log b)$ at the paid width and polynomial-size
weights of the project. This bound assumes the run descriptions are given.
\end{theorem}

\section{How much of the digit survives padding?}
In the detection reduction the primes satisfy $\ell>m\ge2$. Write
$d=\ell-m$ and $\Delta^d g=\sum_{t=0}^d(-1)^{d-t}\binom dt g(t)$.
For any group-algebra element $h$, applying the digit separately
at each of the $d+1$ paddings duplicates work. In fact
\[
\Delta^d D_\ell(h+t[O])
 = \Delta^d\!\left(\frac{t^\ell-t}{\ell}\right)[O]
   +\sum_{j=1}^m \frac{\binom{\ell}{j}}{\ell}
                    (\Delta^d t^{\ell-j})h^j \pmod\ell .
\]
Here the differences are in $t$ and are evaluated at zero. The formula
follows by binomial expansion. The constant $D_\ell(h)$ cancels since
$d\ge1$. If $j>m$ then $\ell-j<d$ and its coefficient also vanishes.
The coefficient of $h^m$ is $d!\binom{\ell}{m}/\ell$, nonzero
modulo $\ell$. Finally let $h=\sum_{i\in S}f_i$ and take the signed
subset sum over $S\subseteq[m]$. Every term of degree below $m$
cancels. The survivor is
\[
   m!d!\,\frac{\binom{\ell}{m}}{\ell}\prod_i f_i
                       =(\ell-1)!\prod_i f_i .
\]
Thus the padded part of the previous reconstruction can be evaluated using
arithmetic modulo $\ell$ and degree at most $m$. This removes an
unnecessary $\ell$th-power calculation for that combined query. The
full mixed coefficient is still present. This calculation alone gives
the original arity cost.

\section{A common group step on independent abscissa masks}
Consider the project's prime-order subclass $r=|G|>2$. Write
$\rho=\tau/2^b$, $\mu=(r-1)\rho$ and let $F_1,\dots,F_m$ be
independently drawn signed hash masks. One affine point is retained
with probability $\rho$. For distinct abscissas joint retention in
one mask has probability at most $\rho^2$, under both the random
odd multiplier and random shift. We use the pair-hash lemma of the
supplied compressed relation-search paper. The following argument
accounts for the identical hash values of $P$ and $-P$.

For nonzero $Q$ define $a_i(Q)=|\{P:P,P+Q\in F_i\}|$. Among
the pairs tested, equal abscissas require $2P=-Q$, which has only
one solution. Therefore $\mathbb E a_i(Q)\le r\rho^2+\rho$.
Independence between slots and the union bound give
\[
 \Pr\big(\exists\,i<j,\ Q\ne O:\ a_i(Q)a_j(Q)>0\big)
       \le \binom{m}{2}(r-1)(r\rho^2+\rho)^2 .                 \tag{1}
\]
At the paid width, for growing polylogarithmic $m$,
$\log_2 r=b+O(1)$ and $\log_2\mu=b/m+O(1)$, since the $2r^{1/m}$
threshold dominates $20m$. The failure bound in (1) is
$2^{-b+4b/m+O(\log m+1)}$.

Outside this event, different masks share no nonzero internal
difference. Consequently for \emph{every} shift $V$ and every
$i\ne j$, $|F_i\cap(F_j+V)|\le1$. Two distinct points in the
intersection would realize a common nonzero difference. For any
step $Q\ne O$ there is also at most one exceptional index $i_*$ with
$a_{i_*}(Q)>0$. These statements hold simultaneously, including for
a step chosen after the seeds are known.

Set $U=\bigcup_i F_i$. Its number of starts of occupied runs along
$Q$ is $R_Q(U)=|U|-|U\cap(U+Q)|$. This counts the boundaries
from empty to occupied in the cyclic ordering generated by $Q$.
Inclusion-exclusion gives $|U|\ge\sum_i|F_i|-\binom m2$.
In the intersection, the exceptional self-pair term is bounded by
$|F_{i_*}|$ and the ordered cross terms by $m(m-1)$. Hence
\[
     R_Q(U)\ge\sum_{i\ne i_*}|F_i|-\tfrac32m(m-1).            \tag{2}
\]
Choose an arbitrary $i_*$ if all self terms vanish. For two masks the
right side is at least $\min(|F_1|,|F_2|)-3$. Adding the padding
point can reduce the number of starts by at most one. The digit
queries have $\ell>m$, so positive weights $1,\ldots,m$ on $U$
are nonzero modulo $\ell$. Even signed run representations require
at least $R_Q$ runs because their group difference has at least $2R_Q$
nonzero jumps, while one run supplies at most two.

\paragraph{A sparse combination of translations.}
The same event limits denominators more general than $1-[Q]$.
With additional probability loss $O(m/\mu)$, all sizes lie in
$[\mu/2,2\mu]$. Indeed, view a slot size as twice a sum over
distinct abscissas. The pair bound gives variance at most $2\mu$,
so Chebyshev and a union bound suffice. Now for every $V\ne O$
the overlap of $U$ with its translate is at most $2\mu+m(m-1)$.
Also $|U|\ge m\mu/2-\binom m2$. Let
$c=\sum_{j=1}^s c_j[V_j]$ have distinct shifts and nonzero
coefficients modulo $\ell$, with $1\le s\le m/16$.
In each translated copy of $U$ there are at least
\[
 m\mu/2-\tbinom m2-(s-1)(2\mu+m(m-1))
\]
positions outside all the other copies. For growing polylogarithmic
$m$ this is at least $m\mu/4$ eventually. At such a position the product
$c\,h$ has exactly one contributing translate, for
$h=\sum_i 1_{F_i}$. Its coefficient is nonzero. Thus
$|\supp(c\,h)|\ge sm\mu/4$, even for data-dependent shifts and
coefficients. Fixed padding changes the union size by at most one
and the overlap estimate by at most two, giving the same asymptotic
bound. This excludes sparse numerators obtained by multiplying the
mixture by a short sum of translations. It does not bound arbitrary
arithmetic circuits for the digit.

\section{Rank of the digit itself at a fixed target}
Here is a separate test of proposals that separate two groups of input
seeds before taking a scalar product. Fix all odd multipliers and $T$.
Let $\ell>m$ be an odd prime distinct from $r$, and split the $m$ slots
into two nonempty groups. For each padding $0\le t\le\ell-m$ let $M_t$ be
the matrix of $D_\ell(t[O]+\sum_i f_i)(T)$, indexed by the left
and right tuples of actual hash shifts. Write $n=(r-1)/2$.
Some allowed $t$ satisfies
\[
 \operatorname{rank}_{\mathbb F_\ell} M_t
   \ \ge\ \frac{n-2\delta-4}{\ell-m+1},\qquad 0\le\delta\le\tau .
\]
In particular the lower bound is $2^{b-O(\log\ell+1)}$ on growing
prime-order inputs at the paid width with $\ell=O(m)$. It is a bound
for an exact separated sum on all shifts. It does not establish
an average runtime lower bound.

For details, $\delta$ is the codimension of the span of shifted slot
masks in the $n$-dimensional space of even functions vanishing at
$O$. The shift-span proof in \texttt{../slot\_arity/rate.tex} applies
in characteristic $\ell$: a vector on abscissas annihilating the
windows, extended by zero to all $2^b$ positions, is $\tau$-periodic
and has sum zero over a period. Its dimension is at most $\tau$.
Differences to one reference shift consequently span at least
$n-\delta-1$ dimensions. Intersecting with the zero-augmentation
space (sum of coefficients zero) leaves at least $n-\delta-2$.

In the space $V$ of all even functions with zero augmentation, of
dimension $n$, the pairing $(f,g)\mapsto(fg)(T)$ is nondegenerate.
After extending scalars the nontrivial Fourier coordinates pair by
inversion, with nonzero target weight $z+z^{-1}$ for an $r$th root
$z$. Vanishing would require $z^2=-1$, impossible as $r$ is odd
and $\ell\ne2$. The Fourier transform is invertible as $\ell\ne r$.
An adjacent nonzero window difference is, up to sign, $e_P$ or
$e_P-e_R$ for $P\ne\pm R$, where $e_P=[P]+[-P]$.
Both act invertibly on $V$. For the difference use
$z+z^{-1}=w+w^{-1}$ only if $z=w$ or $z=w^{-1}$, and faithfulness
of each nontrivial character of the prime-order group.

Keep a free shift on each side and center it by subtracting its
reference. For every extra shift axis take a nonzero adjacent
difference. These are linear row and column operations on $M_t$.
Apply them to the previous signed-subset and padding identity.
Every proper subset omits an axis and hence contributes zero.
Only the full mixture $M_t$ appears for each $t$. On the product
side the centered free slots pair through the extra window
differences, which are invertible on $V$. Restriction to two
subspaces of codimension at most $\delta+2$ gives rank at least
$n-2\delta-4$. The alternating sum over the $\ell-m+1$ compressed
matrices has that rank, proving the bound by subadditivity.

\section{Cost at the target distribution}
Typically (2) forces order $m\mu$ starts for the full mixture. The
class-free formula explicitly visits endpoint populations. A positive
run representation must begin a run at each start, yielding distinct
degree-zero formal atoms, of weights between one and $m$ in the
unpadded full mixture. In particular populations choosing $\ell$
different such atoms have a nonzero scalar weight and pass the degree
cutoff. With $\ell>m$, $\ell=O(m)$, listing these candidate populations
has log cost at least $b-O(m\log m+\log b)$ at these sizes. This is
a cost of listing candidates, not a rank or time lower bound for
methods that sum them implicitly.

Screening seeds for exceptional short descriptions also costs target
coverage. On the event where (1) has no collision, any full mixture
whose support consists of at most $K$ runs for some step satisfies
$\sum_{i\ne i_*}|F_i|\le K+\tfrac32m(m-1)+1$. The extra one allows
padding. Since $|F_{i_*}|\le2\tau$, the total number of targets
reachable by ordinary tuples from this seed is at most
\[
     2\tau\left(\frac{K+\tfrac32m(m-1)+1}{m-1}\right)^{m-1} .
\]
This follows from the product of the slot sizes and the
arithmetic-geometric mean inequality. Divide by $r-1$ for a uniform
nonidentity target independent of the initial seeds, and add (1)
to cover all exceptional seeds. If $m\log_2(K+m^2)=o(b)$, the joint
probability of finding a short description and a reachable target
is at most $2^{-b+o(b)}$. This includes $K=\operatorname{poly}(m)$
and cutoffs with $\log_2 K=O(w/\ell+\log\ell)$ for $w=o(b)$.
In particular selecting only easy seeds in this family would lose
the uniform constant-coverage guarantee. The estimate is for the
specified independent seed distribution and one initial trial.

The best end-to-end bound supplied by the project thus remains
conditional, with log time $O(b/(\log_2 b)^{1/4})$. The missing
saving is a way to evaluate the joint degree-$m$ coefficient on
the actual curve inputs for less than the present arity charge.
The obstruction above is specific to ways of obtaining and
enumerating sparse endpoints. It does not exclude every joint evaluation.

\paragraph{Verification.}
Run \texttt{python -B research/second\_digit/verify.py}.
The evaluator is checked against ordinary integer convolution at
953 targets/parameter choices, including repeated wraps, cancellation,
signed weights, the prime two, and actual elliptic-curve points.
For a length-64 run on the cyclic group of order 1019, prime 11 and
target 346, the digit is 9. Independent integer interval counting
gives 43,077,043,809,387,520 nonconstant orbits. The digit evaluator
does not list them. The hash moment check exhausts 16,384 seed pairs
on the order-19 curve. There are 216 padding checks and 48 direct
digit queries in a reconstruction using three actual masks. At the
actual four-slot paid width 128 on
$y^2=x^3+44x+1$ over $\mathbb F_{65521}$ (order 65,419), two fixed
masks occupy 246 points, with at least 236 runs over every nonzero
step. For four masks the figures are 470 points and 452 runs.
The probability estimate comes from (1), independently of the examples.
\end{document}
